\subsection{Staircase patch geometry}
\label{subsec:peps_torus_staircase_patch_geometry}

\begin{definition}[Horizontal staircase patch and its corner]
    \label{def:peps_horizontal_staircase_patch}
    \lean{TNLean.PEPS.horizontalStaircasePatch}
    \lean{TNLean.PEPS.horizontalStaircaseLeftWindow}
    \lean{TNLean.PEPS.horizontalStaircaseRightWindow}
    \lean{TNLean.PEPS.horizontalStaircaseEndPair}
    \lean{TNLean.PEPS.horizontalStaircaseCorner}
    \lean{TNLean.PEPS.horizontalStaircaseCompletedCorner}
    \leanok
    \uses{def:peps_torus_geometry_translation}
    Fix a horizontal reference edge and write the staircase corner as
    $s=(a,b)$.  For coordinate intervals write
    \[
        I(u,v)=\{t\mid u\le t<v\}.
    \]
    The two end windows of the overlapping-window chain are
    \[
        W_{\mathrm{left}}=I(a,a+L)\times I(b,b+K),
        \qquad
        W_{\mathrm{right}}
        =I(a+L,a+2L)\times I(b+K-1,b+2K-1).
    \]
    Their end pair is $S=W_{\mathrm{left}}\cup W_{\mathrm{right}}$.  The
    staircase patch is the union
    \[
        P=
        \bigl(I(a,a+2L)\times I(b+K-1,b+2K-1)\bigr)
        \cup
        \bigl(I(a,a+L)\times I(b,b+2K-1)\bigr),
    \]
    and the corner cut out by the two end windows is
    \[
        C=I(a,a+L)\times I(b+K,b+2K-1).
    \]
    Completing $C$ by one row gives the $L\times K$ rectangle
    \[
        \widehat C=I(a,a+L)\times I(b+K,b+2K).
    \]
    These cyclic rectangles are the regions used in the staircase chaining and
    corner-stripping step of
    \cite[proof sketch, lines~2320--2445]{Molnar2018NormalPEPS}.
\end{definition}

\begin{lemma}[The staircase end windows are disjoint]
    \label{lem:peps_horizontal_staircase_end_pair_disjoint}
    \lean{TNLean.PEPS.horizontalStaircaseEndPair_disjoint}
    \leanok
    \uses{def:peps_horizontal_staircase_patch}
    If $2L\le W$, then the two end windows of the horizontal staircase are
    disjoint:
    \[
        W_{\mathrm{left}}\cap W_{\mathrm{right}}=\varnothing.
    \]
\end{lemma}
\begin{proof}
    \leanok
    Their horizontal coordinate intervals are $I(a,a+L)$ and $I(a+L,a+2L)$.
    The bound $2L\le W$ prevents wraparound overlap.
\end{proof}

\begin{lemma}[Injectivity of the staircase end windows and completed corner]
    \label{lem:peps_horizontal_staircase_window_injectivity}
    \lean{TNLean.PEPS.NormalTorusArcWindowInjectivityHypotheses.horizontalStaircaseLeftWindow_injective}
    \lean{TNLean.PEPS.NormalTorusArcWindowInjectivityHypotheses.horizontalStaircaseRightWindow_injective}
    \lean{TNLean.PEPS.NormalTorusArcWindowInjectivityHypotheses.horizontalStaircaseCompletedCorner_injective}
    \leanok
    \uses{def:peps_horizontal_staircase_patch}
    Under the one-orientation $L\times K$ torus-window injectivity hypotheses,
    the two end windows and the completed corner are injective regions:
    \[
        W_{\mathrm{left}},\quad W_{\mathrm{right}},\quad \widehat C
        \quad\text{are injective.}
    \]
\end{lemma}
\begin{proof}
    \leanok
    Each of the three regions is an $L\times K$ cyclic rectangle, so the
    window-injectivity hypothesis applies directly.
\end{proof}

\begin{lemma}[Patch decomposition by the end pair and the corner]
    \label{lem:peps_horizontal_staircase_patch_decomposition}
    \lean{TNLean.PEPS.horizontalStaircaseCorner_disjoint_endPair}
    \lean{TNLean.PEPS.horizontalStaircasePatch_eq_endPair_union_corner}
    \lean{TNLean.PEPS.horizontalStaircasePatch_sdiff_endPair}
    \leanok
    \uses{def:peps_horizontal_staircase_patch,
        lem:peps_horizontal_staircase_end_pair_disjoint}
    If $L,K>0$, $2L\le W$, and $2K\le H$, then the corner is disjoint from
    the end pair and the patch decomposes as
    \[
        P=S\cup C,
        \qquad
        P\setminus S=C.
    \]
\end{lemma}
\begin{proof}
    \leanok
    In the no-wraparound rectangle, the rows and columns of the two end
    windows and the corner are the displayed intervals above.  The interval
    calculation gives the disjointness and identifies every point of $P$ with
    either an end-window point or a corner point.
\end{proof}

\begin{theorem}[Fundamental Theorem for normal PEPS on the torus]%
    \label{thm:fundamentalTheorem_normalTorusPEPS}
    \lean{TNLean.PEPS.fundamentalTheorem_normalTorusPEPS_unconditional}
    \lean{TNLean.PEPS.torusAbsorbedGauge_unique_scalar}
    \lean{TNLean.PEPS.torusCovariantAbsorbedGauge_unique_classScalar}
    \lean{TNLean.PEPS.torusGauge_unique_scalar_of_perVertex}
    \leanok
    \uses{thm:fundamentalTheorem_normalSquarePEPS,
        thm:peps_oneVertexComplementComparison,
        def:peps_torus_translation_invariant,
        thm:peps_torusUniformBondDim_of_translationInvariant,
        def:GaugeEquivPEPS}
    Let $A$ and $B$ be translationally invariant normal PEPS on the discrete
    torus such that every $2\times 3$ and $3\times 2$ region is injective. If
    they generate the same state on an $n\times m$ region with $n,m\geq 7$, then
    they are related by horizontal and vertical gauges $X,Y$, up to a scalar
    $\lambda$ satisfying $\lambda^{nm}=1$, and $X$ and $Y$ are unique up to a
    multiplicative constant. This is
    \cite[Theorem~3]{Molnar2018NormalPEPS} on the torus geometry, where every
    edge is interior and translation acts freely.
    Theorem~\ref{thm:fundamentalTheorem_normalTorusPEPS_unconditional} proves
    this at the source's sizes $n,m\geq 7$ with the gauge family described
    by its translation covariance, and the uniqueness of $X$ and $Y$ up to a
    multiplicative constant is proved at the same sizes in
    Theorem~\ref{thm:torusAbsorbedGauge_unique_scalar} per edge,
    Theorem~\ref{thm:torusCovariantAbsorbedGauge_unique_classScalar} on the
    two orientation classes, and
    Theorem~\ref{thm:torusGauge_unique_scalar_of_perVertex} against the
    per-vertex relation. In the torus
    edge convention the two class matrices $X$ and $Y$ are read through this
    translation-covariant family: seam edges carry the transposed inverse of
    the corresponding class matrix, so this is the orientation-faithful form
    of the source's horizontal and vertical gauges.
\end{theorem}

\begin{theorem}[Fundamental Theorem for normal PEPS on the torus,
        translation-covariant form]%
    \label{thm:fundamentalTheorem_normalTorusPEPS_unconditional}
    \lean{TNLean.PEPS.fundamentalTheorem_normalTorusPEPS_unconditional}
    \leanok
    \uses{thm:fundamentalTheorem_normalSquarePEPS,
        thm:peps_oneVertexComplementComparison,
        thm:peps_sameState_transport,
        thm:peps_stateCoeff_translationInvariant,
        thm:peps_torusUniformBondDim_of_translationInvariant,
        def:GaugeEquivPEPS}
    Let $A$ and $B$ be translationally invariant PEPS tensors on the discrete
    $n\times m$ torus with $n,m\geq 7$, with positive physical dimension, equal
    and positive bond dimensions,
    generating the same state, and such that every contiguous $2\times 3$ and
    $3\times 2$ coordinate rectangle of either tensor is injective. Then there
    are a family of invertible matrices, one on each edge, carried onto each
    edge by the torus translations from a single matrix on a reference
    horizontal edge and a single matrix on a reference vertical edge, and a
    single scalar $\lambda$ with $\lambda^{nm}=1$, such that at every site the
    tensor $A$ equals $\lambda$ times the gauge action of the family on $B$.
    The family also realizes the absorbed inserted-coefficient equality at
    every edge: inserting any matrix on an edge of $A$ gives the same
    coefficient as inserting it on the corresponding edge of the gauge-absorbed
    $B$.
    This is \cite[Theorem~3]{Molnar2018NormalPEPS} on the torus at the source's
    sizes $n,m\geq 7$, with the horizontal and vertical gauges $X$, $Y$
    described by translation covariance: the stored orientation of a
    wraparound edge is reversed, so on the seam the family carries the
    transposed inverse of the class matrix, and a family that is literally one
    matrix per orientation class cannot realize the absorbed equality there.
    The translation covariance of the family, with the scalar $\lambda$ kept
    separate rather than absorbed, realizes the remark of
    \cite{Molnar2018NormalPEPS} following the general normal theorem: for a
    translationally invariant system the gauges are also translationally
    invariant, provided the proportionality constants are not absorbed into
    the gauges.
    The uniqueness of $X$ and $Y$ up to a multiplicative constant is not part
    of this statement; the uniqueness clause is
    Theorem~\ref{thm:torusAbsorbedGauge_unique_scalar} (per edge, against the
    absorbed equality), Theorem~\ref{thm:torusCovariantAbsorbedGauge_unique_classScalar}
    (one constant per orientation class), and
    Theorem~\ref{thm:torusGauge_unique_scalar_of_perVertex} (against the
    per-vertex relation itself).
    \begin{proof}
        \leanok
        \uses{lem:peps_injectiveRegion_union_pos}
        The union closure of injective regions follows from the positive bond
        dimensions and Lemma~\ref{lem:peps_injectiveRegion_union_pos}. The
        reference blockings are anchored against the far seam: the removed
        blocks keep a margin of at least two rows and columns below and to
        the left and end exactly at the seam above and to the right, so the
        complementary region decomposes into wraparound-free rectangular
        bands on every torus with $n,m\geq 7$. The
        coherent blocking frames at every edge produce per-edge gauges, whose
        absorption into $B$ gives the absorbed inserted-coefficient equality at
        every edge; the absorbing family is built by transporting one reference
        gauge per orientation class, so it commutes with the torus
        translations. A gauge preserves blocked-region linear independence, so
        the region comparison applies to the gauge-absorbed tensor: comparing
        the corner region with its one-site completion yields proportionality
        scalars $c_R$ and $c_S$ such that, with $\lambda=c_S/c_R$,
        \[
            A_v(\eta,\sigma)
            =
            \lambda\,
            \bigl(\text{gauge action of }X\text{ on }B\bigr)_v(\eta,\sigma).
        \]
        Translation covariance carries the same scalars to every site. Taking
        the product over all vertices and using the invariance of the closed
        state coefficient gives
        \[
            \lambda^{nm}=1.
        \]
    \end{proof}
\end{theorem}

\begin{corollary}[Torus scalar condition from the per-vertex two-block proportionalities]%
    \label{cor:peps_lambda_pow_card_torus_eq_one_of_twoBlockProportional}
    \lean{TNLean.PEPS.lambda_pow_card_torus_eq_one_of_twoBlockProportional}
    \leanok
    \uses{cor:peps_component_eq_gaugeVertex_of_twoBlockProportional}
    Let $A$ and $B$ be PEPS tensors on the discrete $n\times m$ torus with $n,m>1$,
    with equal bond dimensions and positive bond dimensions, generating the same
    state, and let $X$ be a family of invertible matrices, one on each edge. Suppose
    that at every site $v$ a region $R_v$ not containing $v$ is supplied together with
    the two two-block scalar proportionalities of the blocked weights of $A$ and of
    the gauge-absorbed second tensor over $R_v$ and over $R_v\cup\{v\}$, with ratios
    $c_R(v)$ and $c_S(v)$ satisfying $c_R(v)\ne 0$ and whose quotient
    $c_S(v)/c_R(v)$ is a single scalar $\lambda$ independent of $v$, and with the
    gauge-absorbed region block linearly independent. If, in addition, a single
    comparison region whose block and complement block are both blocked-tensor
    injective is supplied, then $\lambda$ raised to the number of torus sites is one:
    \[
        \lambda^{nm}=1.
    \]
\end{corollary}
\begin{proof}\leanok
    At each of the $nm$ sites $v$ the inserted-site scalar extraction of
    Corollary~\ref{cor:peps_component_eq_gaugeVertex_of_twoBlockProportional} turns
    the two two-block proportionalities into the per-vertex gauge relation
    \[
        A_v(\eta,\sigma)
        =
        \lambda\,\widetilde B_v(\eta,\sigma),
    \]
    where $\widetilde B_v$ is the gauge action of $X$ on $B$ at $v$. Contracting
    both sides over the whole torus pulls one factor of $\lambda$ out of each of
    the $nm$ sites, so the closed torus contractions of $A$ and of $\widetilde B$
    satisfy
    \[
        \langle A\rangle_{\mathrm{torus}}
        =
        \lambda^{nm}\,\langle \widetilde B\rangle_{\mathrm{torus}}.
    \]
    The single comparison region and its complement being injective makes this
    common torus contraction nonzero; by hypothesis the quotient
    $c_S(v)/c_R(v)$ is the same scalar $\lambda$ at every site, so
    $\lambda^{nm}=1$.
\end{proof}

\begin{theorem}[Per-edge uniqueness of the torus gauge up to a constant]%
    \label{thm:torusAbsorbedGauge_unique_scalar}
    \lean{TNLean.PEPS.torusAbsorbedGauge_unique_scalar}
    \leanok
    \uses{def:applyGauge, def:peps_edgeInsertedCoeff}
    Let $A$ and $B$ be translationally invariant PEPS tensors on the discrete
    $n\times m$ torus with $n,m\geq 7$, with positive physical dimension,
    equal and positive bond dimensions, generating the same state, and such
    that every contiguous
    $2\times 3$ and $3\times 2$ coordinate rectangle of either tensor is
    injective.  Let $X$ and $X'$ be two families of invertible matrices, one
    on each edge of the torus, and let $e$ be an edge at which both families
    realize the absorbed inserted-coefficient equality: inserting any matrix
    on the edge $e$ of $A$ gives the same coefficient as inserting it on the
    edge $e$ of the gauge-absorbed $B$, for every physical configuration.
    Then the two families are proportional at $e$: $X'_e = c\cdot X_e$ for a
    nonzero scalar $c$.
    This is the uniqueness clause of \cite[Theorem~3]{Molnar2018NormalPEPS}
    read at one edge, at the source's sizes $n,m\geq 7$ as in
    Theorem~\ref{thm:fundamentalTheorem_normalTorusPEPS_unconditional}; the
    source proves it through the determinacy of the conjugator in the
    isomorphism lemma, whose realizing matrix is unique up to a multiplicative
    constant.
    \begin{proof}
        \leanok
        \uses{lem:peps_injectiveRegion_union_pos}
        The rectangle hypotheses produce, along the orientation class of $e$,
        a comparison region having $e$ as its single boundary edge on which
        the second tensor and its complement are injective; the union closure
        of injective regions follows from the positive bond dimensions and
        Lemma~\ref{lem:peps_injectiveRegion_union_pos}.  The bare-edge
        absorbed equality of each family lifts to the region
        inserted-coefficient identity over this region, with the inserted
        matrix conjugated by the orientation-adapted absorption of the
        family's matrix at $e$.  The region-insertion transfer map is
        determined by this identity, so both orientation-adapted absorbing
        matrices induce the same conjugation on every bond matrix:
        \[
            X_{\mathrm{abs}}\,N\,X_{\mathrm{abs}}^{-1}
            =
            X'_{\mathrm{abs}}\,N\,\bigl(X'_{\mathrm{abs}}\bigr)^{-1}
            \qquad\text{for every bond matrix }N.
        \]
        A matrix commuting with the full bond matrix algebra is a scalar
        multiple of the identity, so $X'_{\mathrm{abs}}=c\,X_{\mathrm{abs}}$
        for some $c\ne0$. Unwinding the orientation adaptation gives
        $X'_e=c\cdot X_e$.
    \end{proof}
\end{theorem}

\begin{theorem}[Class constancy of the gauge comparison scalars]%
    \label{thm:torusCovariantAbsorbedGauge_unique_classScalar}
    \lean{TNLean.PEPS.torusCovariantAbsorbedGauge_unique_classScalar}
    \leanok
    \uses{def:applyGauge, def:peps_edgeInsertedCoeff}
    Let $A$, $B$, $X$, $X'$ be as in
    Theorem~\ref{thm:torusAbsorbedGauge_unique_scalar}, with both families
    realizing the absorbed inserted-coefficient equality at every edge, and
    suppose moreover that both families are translation covariant: the matrix
    at any translate of an edge is the matrix at the edge, transposed-inverted
    exactly when the translation swaps the stored endpoint order.  Then there
    are nonzero scalars $c_h$ and $c_v$ such that on every horizontal edge
    stored in the natural cyclic order (the stored second endpoint is the
    cyclic right neighbour of the first) $X'_e = c_h\cdot X_e$, on every
    horizontal edge wrapping the seam (stored with swapped endpoints)
    $X'_e = c_h^{-1}\cdot X_e$, and likewise on the vertical class with $c_v$.
    The inversion on the seam is forced by the ordered-edge convention: the
    covariance carries the transposed inverse of the class matrix onto a seam
    edge, and the transposed inverse of a proportionality inverts its
    constant.  A single constant on the whole class is therefore not the true
    statement in this convention; this is the class form of the uniqueness
    clause of \cite[Theorem~3]{Molnar2018NormalPEPS}.
    \begin{proof}
        \leanok
        \uses{thm:torusAbsorbedGauge_unique_scalar}
        The constants are the per-edge scalars of
        Theorem~\ref{thm:torusAbsorbedGauge_unique_scalar} at the two
        reference edges.  Every edge of a class is the translate of its
        reference edge by a unique translation $\tau$. At an interior edge,
        covariance gives $X_e=\tau\cdot X_{\mathrm{ref}}$, hence
        \[
            X'_e
            =
            \tau\cdot X'_{\mathrm{ref}}
            =
            \tau\cdot(c_h\,X_{\mathrm{ref}})
            =
            c_h\,(\tau\cdot X_{\mathrm{ref}})
            =
            c_h\,X_e.
        \]
        At a seam edge, covariance gives $X_e=(X_{\mathrm{ref}})^{-\top}$,
        hence
        \[
            X'_e
            =
            (X'_{\mathrm{ref}})^{-\top}
            =
            (c_h\,X_{\mathrm{ref}})^{-\top}
            =
            c_h^{-1}(X_{\mathrm{ref}})^{-\top}
            =
            c_h^{-1}X_e.
        \]
        The marker $e_{1,x}+1=e_{2,x}$ distinguishes the two branches, and the
        vertical class is identical with $c_v$ in place of $c_h$.
    \end{proof}
\end{theorem}

\begin{theorem}[Uniqueness of the gauges against the per-vertex relation]%
    \label{thm:torusGauge_unique_scalar_of_perVertex}
    \lean{TNLean.PEPS.torusGauge_unique_scalar_of_perVertex}
    \leanok
    \uses{def:gaugeVertex, def:applyGauge}
    Let $A$ and $B$ be translationally invariant PEPS tensors on the discrete
    $n\times m$ torus with $n,m\geq 7$, with positive physical dimension,
    equal and positive bond dimensions, generating the same state, and such
    that every contiguous
    $2\times 3$ and $3\times 2$ coordinate rectangle of either tensor is
    injective.  Suppose the families $X$ and $X'$ of invertible matrices, one
    on each edge, each realize the per-vertex relation of
    Theorem~\ref{thm:fundamentalTheorem_normalTorusPEPS_unconditional}: at
    every site the tensor $A$ equals $\lambda$ (respectively $\lambda'$)
    times the gauge action of the family on $B$, with $\lambda^{nm} = 1$
    (respectively $\lambda'^{nm} = 1$).  Then the two families are
    proportional at every edge: $X'_e = c\cdot X_e$ for a nonzero scalar $c$
    depending on the edge.
    This is the uniqueness clause of \cite[Theorem~3]{Molnar2018NormalPEPS}
    read against the displayed relation $B = \lambda\cdot(X,Y)A$ itself, at
    the source's sizes $n,m\geq 7$.
    \begin{proof}
        \leanok
        \uses{thm:torusAbsorbedGauge_unique_scalar,
            def:peps_edgeInsertedCoeff}
        The inserted-edge coefficient is multilinear in the site tensors, one
        factor per site, so the per-vertex relation scales every inserted
        coefficient by $\lambda^{nm} = 1$: each family realizes the absorbed
        inserted-coefficient equality at every edge.
        Theorem~\ref{thm:torusAbsorbedGauge_unique_scalar} then gives the
        proportionality at each edge.
    \end{proof}
\end{theorem}
